% !TeX program = xelatex
% Copyright 2006 Konstantin Korikov <lostclus@ua.fm>
%
% This work may be distributed and/or modified under the
% conditions of the LaTeX Project Public License, either version 1.3
% of this license or (at your option) any later version.
% The latest version of this license is in
%   http://www.latex-project.org/lppl.txt
% and version 1.3 or later is part of all distributions of LaTeX
% version 2003/12/01 or later.
%
% This work has the LPPL maintenance status "maintained".
%
% This Current Maintainer of this work is Konstantin Korikov.
%
% This work consists of all files listed in manifest.txt.
%

%
% This is automaticaly generated file, do not edit it.
%
\documentclass[russian,simple]{eskdtext}
\usepackage{xltxtra}
\usepackage{microtype}

\defaultfontfeatures{Scale=MatchLowercase, Ligatures=TeX, Mapping=tex-text}
\setmainfont{Liberation Serif}
\setromanfont{Liberation Serif}
\setmonofont{Liberation Mono}

\newfontfamily{\cyrillicfont}{Liberation Serif} 
\newfontfamily{\cyrillicfontrm}{Liberation Serif}
\newfontfamily{\cyrillicfonttt}{Liberation Mono}
\newfontfamily{\cyrillicfontsf}{Liberation Sans}

\usepackage{minted}
\setminted{
  fontsize=\small,
}

\usepackage{amstext}
\usepackage{amsmath}
\usepackage{calc}
\usepackage{array}
\usepackage{longtable}
\usepackage{graphics}

\title{Коллекция eskdx v0.98}
\ESKDdocName{Руководство пользователя}
\author{Кориков~К.~В.}
\date{2011/05/13}
\renewcommand{\ESKDtheTitleFieldVIIIr}{%
\copyright 2006 Константин Кориков}

\usepackage[unicode]{hyperref}
\hypersetup{pdfkeywords=ЕСКД ГОСТ LaTeX класс пакет}

\newcommand{\gost}[1]{\mbox{ГОСТ}\ \mbox{#1}}
\newcommand{\pkg}[1]{\mbox{\textsf{#1}}}
\newcommand{\cmd}[1]{\texttt{\textbackslash\mbox{#1}}}
\newcommand{\env}[1]{\mbox{\texttt{#1}}}
\newcommand{\opt}[1]{\mbox{\texttt{#1}}}
\newcommand{\optarg}[1]{\mbox{\textit{#1}}}
\newcommand{\val}[1]{\mbox{\texttt{#1}}}
\newenvironment{defs}{\begin{itemize}}{\end{itemize}}
\newcommand{\cmddef}[2]{\item[]\texttt{\textbackslash#1#2}}
\newcommand{\envdef}[3]{\item[]%
  \texttt{\textbackslash begin\{#1\}#2}
  \par\textit{\quad #3}%
  \par\texttt{\textbackslash end\{#1\}}}
\newcommand{\cmdarg}[1]{\mbox{\textrm{\textit{#1}}}}

\begin{document}
\maketitle
\tableofcontents
\newpage
\section{Общие сведения}
\subsection{О коллекции \pkg{eskdx}}
Коллекция пакетов и классов \pkg{eskdx} рассчитана на использование в
среде \LaTeXe\ и предназначена для верстки документов в соответствии с
требованиями Единой системы конструкторской документации.  Коллекция
содержит три класса: \pkg{eskdtext} (для текстовой документации),
\pkg{eskdgraph} (для чертежей и схем) и \pkg{eskdtab} (для документов,
разбитых на графы), а также некоторое число пакетов (некоторые из них
являются внутренними).

\subsection{Возможности коллекции}
Основные возможности \pkg{eskdx}:
\begin{itemize}
\item основная надпись и дополнительные графы по \gost{2.104-2006},
форма 1, 2, 2а, 2б;
\item поддерживаемые форматы листов: А0, А1, А2, А3, А4;
\item поддержка двусторонней печати;
\item настройка элементов текста (абзацы, перечисления, и т.п.) и
рубрикации (разделы, подразделы, пункты, подпункты, приложения) по
\gost{2.105-95};
\item титульный лист и лист утверждения по \gost{2.105-95};
\item спецификация по \gost{2.106-96}, форма 1, 1а, 2, 2а;
\item лист регистрации изменений по \gost{2.503-90};
\item поддержка множества языков (на данный момент русский и
украинский).
\end{itemize}

\newpage
\section{Базовые принципы использования}
Для использования коллекции \pkg{eskdx} необходимо иметь как
минимум базовые знания системы \LaTeXe.

Основа коллекции "--- классы \pkg{eskdtext}, \pkg{eskdgraph} и
\pkg{eskdtab}. Для текстовых документов, содержащих, в основном,
сплошной текст, используют класс \pkg{eskdtext}, для чертежей и схем
"--- \pkg{eskdgraph}, а для документов, содержащих текст, разбитый на
графы "--- \pkg{eskdtab}.

Примечание "--- в большинстве случаев для получения правильного
выходного файла необходимо выполнить как минимум два прохода
интерпретатора.

\subsection{Пример простого документа}
Пример простого текстового документа:
\begin{quote}
\inputminted{latex}{example.tex}
\end{quote}

Опции \opt{russian} и \opt{koi8-r}, на строке 1, задают,
соответственно язык и кодировку документа. Команды на строках 2 по 16
задают информацию о документе, которая используется в титульном листе
и в основной надписи. Строки 19 по 23 "--- тело документа.

\subsection{Опции классов}



\subsubsection{Общие опции всех классов}
Опции, воспринимаемые классами \pkg{eskdtext} и \pkg{eskdgraph}
приведены в таблице~\ref{t:common:opts}.
\begin{longtable}{|%
p{.32\textwidth-2\tabcolsep-1.5\arrayrulewidth}|
p{.68\textwidth-2\tabcolsep-1.5\arrayrulewidth}|}
\caption{Общие опции классов \pkg{eskdtext} и
\pkg{eskdgraph}}\label{t:common:opts}\\
\hline
\multicolumn{1}{|c|}{Опция}&\multicolumn{1}{c|}{Описание}\\\hline
\endfirsthead
\caption*{Продолжение таблицы \ref{t:common:opts}}\\\hline
\multicolumn{1}{|c|}{Опция}&\multicolumn{1}{c|}{Описание}\\\hline
\endhead

\opt{russian}\par\opt{ukrainian}&
Выбирает язык документа, соответственно русский, украинский. По
умолчанию русский.\\\hline
\opt{koi8-r}\par\opt{koi8-u}\par\opt{cp1251}\par\opt{iso8859-5}\par
\opt{cp866}\par\opt{utf8}\par\opt{utf-8}&
Задает кодировку текста, соответственно \mbox{KOI8-R}, \mbox{KOI8-U},
\mbox{CP1251}, \mbox{ISO8859-5}, \mbox{CP866}, \mbox{UTF-8}.\\\hline
\opt{reduceheight=}&
Уменьшает высоту рамки за счет увеличения нижнего поля на
заданную высоту.\\\hline
\opt{emptystyle}&
Включает пустой стиль страниц, т.е. без рамки формы, без основной
надписи и без дополнительных граф.\\\hline
\opt{stitching}\par\opt{nostitching}&
Включает/выключает поле для переплета. По умолчанию поле включено.\\\hline
\opt{columnvii}\par\opt{nocolumnvii}&
Включает/выключает графу 7 основной надписи (номер листа/страницы). По
умолчанию графа 7 включена только для документов, содержащих больше
одного листа\\\hline
\opt{columnviii}\par\opt{nocolumnviii}&
Включает/выключает графу 8 основной надписи (общее количество
листов/страниц). По умолчанию графа 8 включена.\\\hline
\opt{columnsxix}\par\opt{nocolumnsxix}&
Включает/выключает дополнительные графы 19 по 23. По умолчанию
включены.\\\hline
\opt{columnsxxiv}\par\opt{nocolumnsxxiv}&
Включает/выключает дополнительные графы 24 и 25. По умолчанию
выключены.\\\hline
\opt{columnsxxvii}\par\opt{nocolumnsxxvii}&
Включает/выключает дополнительные графы 27 по 30. По умолчанию
выключены.\\\hline
\opt{columnxxxi}\par\opt{nocolumnxxxi}&
Включает/выключает дополнительную графу 31. По умолчанию
включена.\\\hline
\opt{columnxxxii}\par\opt{nocolumnxxxii}&
Включает/выключает дополнительную графу 32 (формат листа). По
умолчанию включена.\\\hline
\opt{simple}&
Тоже самое что \opt{nocolumnsxix}, \opt{nocolumnxxxi} и
\opt{nocolumnxxxii}.\\\hline
\opt{linethick=}&
Устанавливает толщину сплошной толстой линии. По умолчанию
\val{.4mm}.\\\hline
\opt{linethin=}&
Устанавливает толщину сплошной тонкой линии. По умолчанию
\val{.2mm}.\\\hline
\opt{hpadding=}&
Устанавливает расстояние от рамки формы до границ текста в начале и в
конце строк. По умолчанию \val{3mm}.\\\hline
\opt{vpadding=}&
Устанавливает расстояние от верхней или нижней строки текста до
верхней или нижней рамки. По умолчанию \val{10mm}.\\\hline
\end{longtable}

Примечание "--- Схемы расположения граф основной надписи приведены в
приложении~\ref{a:columnlayout}.

\subsubsection{Опции класса \pkg{eskdtext}}
Опции, воспринимаемые классом \pkg{eskdtext} приведены в
таблице~\ref{t:eskdtext:opts}.
\begin{longtable}{|%
p{.32\textwidth-2\tabcolsep-1.5\arrayrulewidth}|
p{.68\textwidth-2\tabcolsep-1.5\arrayrulewidth}|}
\caption{Опции класса \pkg{eskdtext}}\label{t:eskdtext:opts}\\
\hline
\multicolumn{1}{|c|}{Опция}&\multicolumn{1}{c|}{Описание}\\\hline
\endfirsthead
\caption*{Продолжение таблицы \ref{t:eskdtext:opts}}\\\hline
\multicolumn{1}{|c|}{Опция}&\multicolumn{1}{c|}{Описание}\\\hline
\endhead

\opt{12pt}\par\opt{14pt}&
Устанавливает размер шрифта документа, соответственно \val{12pt},
\val{14pt}. По умолчанию \val{14pt}.\\\hline
\opt{pointsingle}\par\opt{pointsection}\par\opt{pointsubsection}&
Устанавливает стиль нумерации пунктов, соответственно в пределах всего
документа, в пределах раздела, в пределах подраздела. По умолчанию
пункты нумеруются в пределах всего документа (сквозная
нумерация).\\\hline
\opt{floatsingle}\par\opt{floatsection}\par\opt{floatsubsection}&
Устанавливает стиль нумерации плавающих объектов (рисунки, таблицы),
соответственно в пределах всего документа, в пределах раздела, в
пределах подраздела. По умолчанию плавающие объекты нумеруются в
пределах всего документа (сквозная нумерация).\\\hline
\opt{equationsingle}\par\opt{equationsection}\par\opt{equationsubsection}&
Устанавливает стиль нумерации формул, соответственно в пределах всего
документа, в пределах раздела, в пределах подраздела. По умолчанию
формулы нумеруются в пределах всего документа (сквозная
нумерация).\\\hline
\opt{footnotearabic}\par\opt{footnoteasterisk}&
Устанавливает стиль нумерации сносок, соответственно
арабской цифрой с закрывающей круглой скобкой,
либо звездочками (количество звездочек соответствует
номеру сноски). Нумерация осуществляется в пределах 
страницы. По умолчанию используется \val{footnotearabic}\\\hline
\opt{twoside}&
Включает двустороннюю печать.\\\hline
\opt{draft}&
Включает черновой режим.\\\hline
\end{longtable}

\subsubsection{Опции класса \pkg{eskdgraph}}
Опции, воспринимаемые классом \pkg{eskdgraph} приведены в
таблице~\ref{t:eskdgraph:opts}.
\begin{longtable}{|%
p{.32\textwidth-2\tabcolsep-1.5\arrayrulewidth}|
p{.68\textwidth-2\tabcolsep-1.5\arrayrulewidth}|}
\caption{Опции класса \pkg{eskdgraph}}\label{t:eskdgraph:opts}\\
\hline
\multicolumn{1}{|c|}{Опция}&\multicolumn{1}{c|}{Описание}\\\hline
\endfirsthead
\caption*{Продолжение таблицы \ref{t:eskdgraph:opts}}\\\hline
\multicolumn{1}{|c|}{Опция}&\multicolumn{1}{c|}{Описание}\\\hline
\endhead

\opt{a0paper}\par\opt{a1paper}\par\opt{a2paper}\par\opt{a3paper}\par
\opt{a4paper}&
Устанавливает формат (размер) листа, соответственно А0, А1, А2, А3,
А4. По умолчанию А4.\\\hline
\opt{landscape}\par\opt{portrait}&
Устанавливает ориентацию страницы, соответственно ландшафтная
ориентация (основная надпись вдоль длинной стороны листа), портретная
ориентация (основная надпись вдоль короткой стороны листа). По
умолчанию ландшафтная ориентация для листов формата больше А4, и
портретная ориентация для листов формата А4.\\\hline
\opt{ignorestamp}\par\opt{noignorestamp}&
Включает/выключает игнорирование основной надписи при
позиционировании рисунка чертежа. По умолчанию основная надпись
игнорируется.\\\hline
\opt{columnxxvi}\par\opt{nocolumnxxvi}&
Включает/выключает дополнительную графу 26. По умолчанию графа
включена.\\\hline
\opt{columnxxxiii}\par\opt{nocolumnxxxiii}&
Включает/выключает дополнительную графу 33. По умолчанию
выключена.\\\hline
\opt{columnxxxiv}\par\opt{nocolumnxxxiv}&
Включает/выключает дополнительную графу 34. По умолчанию
выключена.\\\hline
\opt{zonelabels}\par\opt{nozonelabels}&
Включает/выключает отметки зон. По умолчанию отметки выключены.\\\hline
\end{longtable}

\subsubsection{Опции класса \pkg{eskdtab}}
Опции, воспринимаемые классом \pkg{eskdtab} приведены в
таблице~\ref{t:eskdtab:opts}.
\begin{longtable}{|%
p{.32\textwidth-2\tabcolsep-1.5\arrayrulewidth}|
p{.68\textwidth-2\tabcolsep-1.5\arrayrulewidth}|}
\caption{Опции класса \pkg{eskdtab}}\label{t:eskdtab:opts}\\
\hline
\multicolumn{1}{|c|}{Опция}&\multicolumn{1}{c|}{Описание}\\\hline
\endfirsthead
\caption*{Продолжение таблицы \ref{t:eskdtab:opts}}\\\hline
\multicolumn{1}{|c|}{Опция}&\multicolumn{1}{c|}{Описание}\\\hline
\endhead

\opt{a3paper}\par\opt{a4paper}&
Устанавливает формат (размер) листа, соответственно А3, А4. По
умолчанию А4.\\\hline
\opt{landscape}\par\opt{portrait}&
Устанавливает ориентацию страницы, соответственно ландшафтная
ориентация (основная надпись вдоль длинной стороны листа), портретная
ориентация (основная надпись вдоль короткой стороны листа). По
умолчанию ландшафтная ориентация для листов формата А3, и
портретная ориентация для листов формата А4.\\\hline
\opt{twoside}&
Включает двустороннюю печать.\\\hline
\opt{draft}&
Включает черновой режим.\\\hline
\end{longtable}

\subsection{Информация о документе}
\label{docinfo}
Информацию о документе размещают в преамбуле документа используя
описанные ниже команды.

Наименование ведомства, в систему которого входит организация,
разработавшая документ, задает команда
\begin{defs}
\cmddef{ESKDdepartment}{\{\cmdarg{текст}\}}
\end{defs}

Наименование организации задает команда
\begin{defs}
\cmddef{ESKDcompany}{\{\cmdarg{текст}\}}
\end{defs}

Код по классификатору продукции задает команда
\begin{defs}
\cmddef{ESKDclassCode}{\{\cmdarg{текст}\}}
\end{defs}

Наименование изделия задают команды
\begin{defs}
\cmddef{ESKDunitName}{\{\cmdarg{текст}\}}
\cmddef{ESKDtitle}{\{\cmdarg{текст}\}}
\cmddef{title}{\{\cmdarg{текст}\}}
\end{defs}
все три команды действуют одинаково (команды близнецы).

Наименование документа задает команда
\begin{defs}
\cmddef{ESKDdocName}{\{\cmdarg{текст}\}}
\end{defs}

Обозначение документа задает команда
\begin{defs}
\cmddef{ESKDsignature}{\{\cmdarg{текст}\}}
\end{defs}

Литеру, присвоенную данному документу, задает команда
\begin{defs}
\cmddef{ESKDletter}{\{\cmdarg{буква1}\}\{\cmdarg{буква2}\}%
\{\cmdarg{буква3}\}}
\end{defs}

Обозначение материала детали задает команда
\begin{defs}
\cmddef{ESKDmaterial}{\{\cmdarg{текст}\}}
\end{defs}

Массу изделия задает команда
\begin{defs}
\cmddef{ESKDmass}{\{\cmdarg{текст}\}}
\end{defs}

Масштаб чертежа задает команда
\begin{defs}
\cmddef{ESKDscale}{\{\cmdarg{текст}\}}
\end{defs}

Наименование или различительный индекс предприятия задает команда
\begin{defs}
\cmddef{ESKDgroup}{\{\cmdarg{текст}\}}
\end{defs}

Фамилию автора задают команды
\begin{defs}
\cmddef{ESKDauthor}{\{\cmdarg{текст}\}}
\cmddef{author}{\{\cmdarg{текст}\}}
\end{defs}

Фамилию проверяющего задает команда
\begin{defs}
\cmddef{ESKDchecker}{\{\cmdarg{текст}\}}
\end{defs}

Фамилию нормоконтролера задает команда
\begin{defs}
\cmddef{ESKDnormContr}{\{\cmdarg{текст}\}}
\end{defs}

Фамилию утвердившего задает команда
\begin{defs}
\cmddef{ESKDapprovedBy}{\{\cmdarg{текст}\}}
\end{defs}

Дату документа задают команды
\begin{defs}
\cmddef{ESKDdate}{\{\cmdarg{год/месяц/число}\}}
\cmddef{date}{\{\cmdarg{год/месяц/число}\}}
\end{defs}

Описанные выше команды сохраняют значения своих аргументов в
соответствующие \cmd{ESKDthe}-команды. Например, команда
\cmd{ESKDtheAuthor} печатает фамилию автора, а команда
\cmd{ESKDtheNormContr} печатает фамилию нормоконтролера.

\sloppy
Немного иначе обстоят дела с датой и литерой. Год печатает команда
\cmd{ESKDtheYear}, месяц "--- \cmd{ESKDtheMonth}, а число "---
\cmd{ESKDtheDay}. Буквы литеры печатают команды \cmd{ESKDtheLetterI},
\cmd{ESKDtheLetterII}, \cmd{ESKDtheLetterIII}.

\fussy

\subsection{Титульный лист}
Титульный лист печатает команда
\begin{defs}
\cmddef{maketitle}{}
\end{defs}
которую размещают в теле документа.

Согласно \gost{2.105} титульный лист разбит на 10 полей. По
умолчанию большинство полей уже содержат необходимые данные, заданные
командами, описанными в разделе~\ref{docinfo}. Также для заполнения
титульного листа существует несколько дополнительных команд, которые
размещают в преамбуле документа.

Команда
\begin{defs}
\cmddef{ESKDtitleApprovedBy}{\{\cmdarg{должность}\}\{\cmdarg{фамилия}\}}
\end{defs}
размещает гриф утверждения в правой части поля 3.

Команда
\begin{defs}
\cmddef{ESKDtitleAgreedBy}{\{\cmdarg{должность}\}\{\cmdarg{фамилия}\}}
\end{defs}
размещает гриф согласования в левой части поля 3.  Если согласовавших
лиц больше одного, то команду повторяют несколько раз, при этом
остальные фамилии размещаются в левой части поля 8.

Список разработчиков документа задают при помощи команды
\begin{defs}
\cmddef{ESKDtitleDesignedBy}{\{\cmdarg{должность}\}\{\cmdarg{фамилия}\}}
\end{defs}
Команду повторяют для каждого разработчика. Список размещается в
правой части поля 8.

При утверждении документа листом утверждения его обозначение указывают
в левой части поля 3. Для этого используют команду
\begin{defs}
\cmddef{ESKDtitleApprovingSheet}{\{\cmdarg{обозначение}\}}
\end{defs}

Кроме описанных выше методов заполнения титульного листа, имеется
метод, позволяющий принудительного указать какие данные нужно
разместить в том или ином поле посредством переопределения
\cmd{ESKDtheTitleField}-команд. Например, для того чтобы в поле 10
разместить какой-нибудь текст (по умолчанию поле содержит год) в
преамбулу документа нужно поместить
\begin{quote}
\begin{minted}{latex}
\renewcommand{\ESKDtheTitleFieldX}{%
какой-нибудь текст}
\end{minted}
\end{quote}

Примечание "--- Схема расположения полей титульного листа и листа
утверждения приведена в приложении~\ref{a:titlelayout}.

\subsection{Заполнение граф основной надписи и дополнительных граф}
По умолчанию большинство граф основной надписи и некоторые
дополнительные графы содержат данные, заданные командами, описанными в
разделе~\ref{docinfo}. Для принудительного заполнения граф
предназначены \cmd{ESKDcolumn}-команды, которые размещают в преамбуле
документа.

Графу 1 (наименование изделия/документа) задает команда
\begin{defs}
\cmddef{ESKDcolumnI}{\{\cmdarg{текст}\}}
\end{defs}

Графу 2 (обозначение документа) задает команда
\begin{defs}
\cmddef{ESKDcolumnII}{\{\cmdarg{текст}\}}
\end{defs}

Графу 3 (обозначение материала детали) задает команда
\begin{defs}
\cmddef{ESKDcolumnIII}{\{\cmdarg{текст}\}}
\end{defs}

Графу 4 (литера) задают команды
\begin{defs}
\cmddef{ESKDcolumnIVfI}{\{\cmdarg{текст}\}}
\cmddef{ESKDcolumnIVfII}{\{\cmdarg{текст}\}}
\cmddef{ESKDcolumnIVfIII}{\{\cmdarg{текст}\}}
\end{defs}

Графу 5 (масса изделия) задает команда
\begin{defs}
\cmddef{ESKDcolumnV}{\{\cmdarg{текст}\}}
\end{defs}

Графу 6 (масштаб) задает команда
\begin{defs}
\cmddef{ESKDcolumnVI}{\{\cmdarg{текст}\}}
\end{defs}

Графу 9 (наименование или различительный индекс предприятия) задает
команда
\begin{defs}
\cmddef{ESKDcolumnIX}{\{\cmdarg{текст}\}}
\end{defs}

Графу 11 (фамилии лиц, подписывающих документ) задают команды
\begin{defs}
\cmddef{ESKDcolumnXIfI}{\{\cmdarg{текст}\}}
\cmddef{ESKDcolumnXIfII}{\{\cmdarg{текст}\}}
\cmddef{ESKDcolumnXIfIII}{\{\cmdarg{текст}\}}
\cmddef{ESKDcolumnXIfIV}{\{\cmdarg{текст}\}}
\cmddef{ESKDcolumnXIfV}{\{\cmdarg{текст}\}}
\cmddef{ESKDcolumnXIfVI}{\{\cmdarg{текст}\}}
\end{defs}

Графу 19 (инвентарный номер подлинника) задает команда
\begin{defs}
\cmddef{ESKDcolumnXIX}{\{\cmdarg{текст}\}}
\end{defs}

Графу 21 (инвентарный номер подлинника, взамен которого выпущен данный
подлинник) задает команда
\begin{defs}
\cmddef{ESKDcolumnXXI}{\{\cmdarg{текст}\}}
\end{defs}

Графу 22 (инвентарный номер дубликата) задает команда
\begin{defs}
\cmddef{ESKDcolumnXXII}{\{\cmdarg{текст}\}}
\end{defs}

Графу 24 (обозначение документа, взамен или на основании которого
выпущен документ) задает команда
\begin{defs}
\cmddef{ESKDcolumnXXIV}{\{\cmdarg{текст}\}}
\end{defs}

Графу 25 (обозначение соответствующего документа, в котором впервые
записан документ) задает команда
\begin{defs}
\cmddef{ESKDcolumnXXV}{\{\cmdarg{текст}\}}
\end{defs}

Графу 27 (знак, установленный заказчиком) задает команда
\begin{defs}
\cmddef{ESKDcolumnXXVII}{\{\cmdarg{текст}\}}
\end{defs}

Графу 28 (номер решения и год утверждения документации соответствующей
литеры) задает команда
\begin{defs}
\cmddef{ESKDcolumnXXVIII}{\{\cmdarg{текст}\}}
\end{defs}

Графу 29 (номер решения и год утверждения документации) задает команда
\begin{defs}
\cmddef{ESKDcolumnXXIX}{\{\cmdarg{текст}\}}
\end{defs}

Графу 30 (индекс заказчика) задает команда
\begin{defs}
\cmddef{ESKDcolumnXXX}{\{\cmdarg{текст}\}}
\end{defs}

Графу 32 (обозначение формата листа) задает команда
\begin{defs}
\cmddef{ESKDcolumnXXXII}{\{\cmdarg{текст}\}}
\end{defs}

Графу 34 (номера авторских свидетельств на изобретения) задает команда
\begin{defs}
\cmddef{ESKDcolumnXXXIV}{\{\cmdarg{текст}\}}
\end{defs}

Описанные выше команды сохраняют значения своих аргументов в
соответствующие \cmd{ESKDthe}-команды. Например, команда
\cmd{ESKDtheColumnI} печатает содержимое графы 1.

Примечание "--- Схемы расположения граф основной надписи приведены в
приложении~\ref{a:columnlayout}.

\subsection{Рубрикация}
Разделы, подразделы, подподразделы, пункты, подпункты, подподпункты
создаются соответственно командами:
\begin{defs}
\cmddef{section}{\{\cmdarg{заголовок}\}}
\cmddef{subsection}{\{\cmdarg{заголовок}\}}
\cmddef{subsubsection}{\{\cmdarg{заголовок}\}}
\cmddef{point}{}
\cmddef{subpoint}{}
\cmddef{subsubpoint}{}
\end{defs}

По умолчанию пункты нумеруются в пределах всего документа. Чтобы
включить нумерацию пунктов в пределах раздела или подраздела
используют опцию \opt{pointsection} или \opt{pointsubsection} класса
\pkg{eskdtext} соответственно.

Приложение создает команда
\begin{defs}
\cmddef{ESKDappendix}{\{\cmdarg{тип}\}\{\cmdarg{заголовок}\}}
\end{defs}
Первый аргумент "--- тип приложения (обязательное, рекомендуемое или
справочное). Второй аргумент "--- заголовок приложения.

\subsection{Пояснения символов, входящих в формулу}
\sloppy
Командные скобки
\begin{defs}
\envdef{ESKDexplanation}{[ширина]}{пояснения}
\end{defs}
создают список пояснений символов. Перед первым пояснением ставят
\cmd{item[где ]}. Перед последующими пояснениями ставят \cmd{item}.
Окружение может принимать один необязательный параметр~--- ширину
маркера \mbox{<<где >>} в единицах, понятных \LaTeX. Это может быть
полезно в случаях, когда значение по умолчанию не подходит, либо вы
используете другой маркер вместо \mbox{<<где >>}. Для автоматического
вычисления размера маркера можно воспользоваться командой
\cmd{settowidth}.

Длина \cmd{ESKDexplanIndent} задает отступ списка пояснений от левого
края страницы. По умолчанию \val{0pt}.

\fussy

\subsection{Лист регистрации изменений}
Пакет \pkg{eskdchngsheet} предоставляет командные скобки
\begin{defs}
\envdef{ESKDchangeSheet}{}{строки}
\end{defs}
которые создают лист регистрации изменений. По сути это таблица из
десяти колонок, которую заполняют как обычную таблицу, разделяя данные
ячеек символом \verb|&|, а строки "--- \verb|\\|.

\subsection{Чертежи и схемы}
Для чертежей и схем используют класс \pkg{eskdgraph}. Команды включения
графики размещают внутри командных скобок
\begin{defs}
\envdef{ESKDdrawing}{}{графика}
\end{defs}
Для каждой страницы используют свое окружение \env{ESKDdrawing}.

\subsection{Спецификация}
Для спецификации используют класс \pkg{eskdtab} и пакет
\pkg{eskdspec}. Создают спецификацию командные скобки
\begin{defs}
\envdef{ESKDspecification}{}{строки}
\end{defs}
По сути это таблица из семи колонок, которую заполняют как обычную
таблицу, разделяя данные ячеек символом \verb|&|, а строки "---
\verb|\\|.

\subsection{Спецификация при плазовом методе}
Для спецификации при плазовом методе (форма 2 и 2а) используют класс
\pkg{eskdtab} и пакет \pkg{eskdspecii}. Создают спецификацию командные
скобки
\begin{defs}
\envdef{ESKDspecificationII}{}{строки}
\end{defs}
По сути это таблица из девяти колонок, которую заполняют как обычную
таблицу, разделяя данные ячеек символом \verb|&|, а строки "---
\verb|\\|.

\subsection{Лист утверждения}
Лист утверждения создают как отдельный документ с подключением пакета
\pkg{eskdappsheet}. Заполняют лист утверждения таким же образом как и
титульный лист.

Команда
\begin{defs}
\cmddef{ESKDapprovingDoc}{\{\cmdarg{обозначение}\}}
\end{defs}
размещает в нижней части поля 8 обозначение документа, на который
распространяется данный лист утверждения. Если таких документов
несколько, команду повторяют для каждого документа.

Печатает лист утверждения команда
\begin{defs}
\cmddef{maketitle}{}
\end{defs}
или команда
\begin{defs}
\cmddef{ESKDmakeApprovingSheet}{}
\end{defs}

\subsection{Количество рисунков, таблиц, приложений, и т.д.}
Пакет \pkg{eskdtotal} определяет команду
\begin{defs}
\cmddef{ESKDtotal}{\{\cmdarg{счетчик}\}}
\end{defs}
при помощи которой в текст можно подставлять общее число страниц,
рисунков, таблиц, приложений и литературных источников. Допустимые
значения аргумента \cmdarg{счетчик} соответственно: \val{page},
\val{figure}, \val{table}, \val{appendix} и \val{bibitem}. Например,
команда \cmd{ESKDtotal\{figure\}} печатает количество рисунков.

\newpage
\section{Тонкая настройка}
\subsection{Управление стилями страниц}
\label{s:pagestyles}
Каждая страница оформляется в соответствии с определенным стилем
страницы. Существуют такие предопределенные стили страниц:
\val{empty}, \val{title}, \val{formI}, \val{formII}, \val{formIIab},
\val{freesize}.
Стиль \val{empty} "--- это пустой стиль, в нем нет ни рамки, ни
основной надписи, ни дополнительных граф. Стиль \val{title}
предназначен для титульного листа, в нем есть рамка и дополнительные
графы 19 по 23. Стили \val{formI} и \val{formII} соответствуют форме~1
и форме~2 по \gost{2.104} соответственно, имеют рамку, основную
надпись и дополнительные графы.  Стиль \val{formIIab} соответствует
форме 2б на четной странице при двусторонней печати, и форме 2а в
остальных случаях. Стиль \val{freesize} используется для листов
произвольного размера, подробное описание приведено в 
разделе~\ref{s:freesize}

Менять стиль страницы можно используя команды
\begin{defs}
\cmddef{ESKDstyle}{\{стиль\}}
\cmddef{ESKDthisStyle}{\{стиль\}}
\end{defs}
Первая меняет стиль текущей и последующих страниц, вторая меняет стиль
только текущей страницы.

Также можно изменять существующие стили и добавлять собственные.

Новый стиль создает команда
\begin{defs}
\cmddef{ESKDnewStyle}{\{\cmdarg{стиль}\}\{\cmdarg{высота}\}}
\end{defs}
Первый аргумент задает имя стиля, второй "--- высоту основной надписи.

Команда
\begin{defs}
\cmddef{ESKDputOnStyle}{\{\cmdarg{стиль}\}\{\cmdarg{элемент}\}%
\{\cmdarg{команды}\}}
\end{defs}
заменяет или добавляет в существующий \cmdarg{стиль} \cmdarg{элемент},
для верстки которого используются \cmdarg{команды}. Первый аргумент
"--- имя существующего стиля. Второй аргумент "--- имя элемента,
который нужно добавить или заменить. Третий аргумент "--- команды,
рисующие данный элемент (допустимы только команды окружения
\env{picture}).

Для удаления элементов из стиля используется команда
\begin{defs}
\cmddef{ESKDremoveFromStyle}{\{\cmdarg{стиль}\}\{\cmdarg{элемент}\}}
\end{defs}

В предопределенных стилях рамку формы рисует элемент \val{frame}, а
основную надпись "--- \val{stamp}. Например, для того чтобы удалить
основную надпись в стиле \val{formI} можно воспользоваться командой
\begin{quote}
\verb|\ESKDremoveFromStyle{formI}{stamp}|
\end{quote}

Команды
\begin{defs}
\cmddef{ESKDdefaultTitleStyle}{\{\cmdarg{стиль}\}}
\cmddef{ESKDdefaultFirstStyle}{\{\cmdarg{стиль}\}}
\cmddef{ESKDdefaultStyle}{\{\cmdarg{стиль}\}}
\end{defs}
задают стиль по умолчанию для титульного листа, первой страницы (не
считая титульный лист) и остальных страниц соответственно. Нет смысла
использовать эти команды вне преамбулы документа.

\subsection{Настройка шрифтов}
Команды
\begin{defs}
\cmddef{ESKDfontII}{}
\cmddef{ESKDfontIII}{}
\cmddef{ESKDfontV}{}
\cmddef{ESKDfontVII}{}
\cmddef{ESKDfontX}{}
\end{defs}
определяют базовые шрифты, соответственно размером $2{,}5$, $3{,}5$,
$5$, $7$ и $10~\text{мм}$. Шрифты, заданные этими командами,
используются в основной надписи и в дополнительных графах, а также по
умолчанию в других элементах, таких как титульный лист, спецификация и
лист регистрации изменений.

По умолчанию эти команды устанавливают
\cmd{baselinestretch} в
\begin{defs}
\cmddef{ESKDfontBaseLineStretch}{}
\end{defs}
Устанавливают размер используя соответственно команды 
\begin{defs}
\cmddef{ESKDfontIIsize}{}
\cmddef{ESKDfontIIIsize}{}
\cmddef{ESKDfontVsize}{}
\cmddef{ESKDfontVIIsize}{}
\cmddef{ESKDfontXsize}{}
\end{defs}
Выполняют \cmd{selectfont} и устанавливают начертание используя
команду
\begin{defs}
\cmddef{ESKDfontShape}{}
\end{defs}

Шрифты элементов, представленных в виде таблицы, определяют команды
\begin{defs}
\cmddef{ESKDfontTabHead}{}
\cmddef{ESKDfontTabBody}{}
\end{defs}
Первая определяет шрифт заголовка, вторая "--- шрифт тела таблицы. По
умолчанию определены как \cmd{ESKDfontIII}.

Все приведенные выше команды можно переопределять при помощи
\cmd{renewcommand}.

\subsection{Настройка титульного листа}
Команды
\begin{defs}
\cmddef{ESKDtitleFontI}{}
\cmddef{ESKDtitleFontII}{}
\cmddef{ESKDtitleFontIII}{}
\cmddef{ESKDtitleFontIV}{}
\cmddef{ESKDtitleFontV}{}
\cmddef{ESKDtitleFontVI}{}
\cmddef{ESKDtitleFontVII}{}
\cmddef{ESKDtitleFontVIII}{}
\cmddef{ESKDtitleFontX}{}
\end{defs}
определяют шрифты полей титульного листа. Их можно переопределить
посредством \cmd{renewcommand}.

\subsection{Управление заголовками рубрикации}
Заголовки разделов, подразделов и подподразделов можно выравнивать
влево, вправо, по центру, а также можно задавать стиль отображения и
вертикальные отступы.

Команда
\begin{defs}
\cmddef{ESKDsectAlign}{\{\cmdarg{элемент}\}\{\cmdarg{выравнивание}\}}
\end{defs}
устанавливает выравнивание элемента. Первый аргумент "--- одно
из значений: \val{section}, \val{subsection}, \val{subsubsection};
соответственно для разделов, подразделов, подподразделов. Второй
аргумент "--- одно из значений: \val{Left}, \val{Right}, \val{Center};
соответственно для выравнивания по левому краю, по правому краю, по
центру.

\sloppy
Стиль, в котором будут верстаться заголовки, задает команда
\begin{defs}
\cmddef{ESKDsectStyle}{\{\cmdarg{элемент}\}\{\cmdarg{стиль}\}}
\end{defs}
Первый аргумент "--- одно из значений: \val{section},
\val{subsection}, \val{subsubsection}; соответственно для разделов,
подразделов, подподразделов. Второй аргумент "--- стиль, который может
содержать в себе такие команды как \cmd{Large}, \cmd{bfseries}, и т.п.

Вертикальные отступы задает команда
\begin{defs}
\cmddef{ESKDsectSkip}{\{\cmdarg{элемент}\}\{\cmdarg{отступ перед}\}%
\{\cmdarg{отступ после}\}}
\end{defs}
Первый аргумент "--- одно из значений: \val{section},
\val{subsection}, \val{subsubsection}; соответственно для разделов,
подразделов, подподразделов. Второй аргумент "--- отступ перед
заголовком. Третий аргумент "--- отступ после заголовка.

\fussy

\subsection{Настройка листа регистрации изменений}
\sloppy
Шрифт заголовка и тела таблицы определяют команды
\cmd{ESKDfontTabHead} и \cmd{ESKDfontTabBody}, их можно переопределить
посредством \cmd{renewcommand}.

Примечание "--- команды \cmd{ESKDchngFontHead} и
\cmd{ESKDchngFontBody} являются устаревшими и будут удалены в будущих
версиях \pkg{eskdx}.

\fussy

Ширину колонок таблицы определяют команды
\begin{defs}
\cmddef{ESKDchngColumnIwidth}{}
\cmddef{ESKDchngColumnIIwidth}{}
\cmddef{ESKDchngColumnIIIwidth}{}
\cmddef{ESKDchngColumnIVwidth}{}
\cmddef{ESKDchngColumnVwidth}{}
\cmddef{ESKDchngColumnVIwidth}{}
\cmddef{ESKDchngColumnVIIwidth}{}
\cmddef{ESKDchngColumnVIIIwidth}{}
\cmddef{ESKDchngColumnIXwidth}{}
\cmddef{ESKDchngColumnXwidth}{}
\end{defs}
Их также можно переопределить посредством \cmd{renewcommand}.

\subsection{Настройка спецификации}
\sloppy
Шрифт заголовка и тела таблицы определяют команды
\cmd{ESKDfontTabHead} и \cmd{ESKDfontTabBody}, их можно переопределить
посредством \cmd{renewcommand}.

Примечание "--- команды \cmd{ESKDspecFontHead} и
\cmd{ESKDspecFontBody} являются устаревшими и будут удалены в будущих
версиях \pkg{eskdx}.

\fussy

\subsection{Текстовые документы без рамок формы и основной надписи}
Опция \opt{emptystyle} класса \pkg{eskdtext} позволяет получить
документ без таких элементов как рамка формы, основная надпись и
дополнительные графы. По сути эта опция устанавливает для всех страниц
стиль \val{empty}. Этот стиль не рекомендуется модифицировать, так что
если необходимо добавить на страницу произвольные элементы, нужно
создать и установить новый стиль, как это описано в
разделе~\ref{s:pagestyles}.

Пакет \pkg{eskdplain} создает новый стиль страницы \val{plain},
который не содержит никаких элементов, кроме элемента
\val{pagenumber}, который печатает по центру в нижней части страницы ее
номер. Расположение номера на странице можно контролировать опциями
пакета, которые приведены в таблице~\ref{t:eskdplain:opts}.

\begin{longtable}{|%
p{.32\textwidth-2\tabcolsep-1.5\arrayrulewidth}|
p{.68\textwidth-2\tabcolsep-1.5\arrayrulewidth}|}
\caption{Опции пакета \pkg{eskdplain}}\label{t:eskdplain:opts}\\
\hline
\multicolumn{1}{|c|}{Опция}&\multicolumn{1}{c|}{Описание}\\\hline
\endfirsthead
\caption*{Продолжение таблицы \ref{t:eskdplain:opts}}\\\hline
\multicolumn{1}{|c|}{Опция}&\multicolumn{1}{c|}{Описание}\\\hline
\endhead

\opt{numbertop}\par\opt{numberbottom}&
Располагает номер в верхней или в нижней части страницы
соответственно.\\\hline
\opt{numberleft}\par\opt{numbercenter}\par\opt{numberright}&
При \emph{односторонней} компоновке документа располагает номер 
слева, по центру или справа соответственно.\par
При \emph{двусторонней} компоновке располагает
номер у переплета, по центру или
у внешнего края страниц соответственно.\\\hline
\end{longtable}

Пакет \pkg{eskdplain} устанавливает стиль \val{empty} в качестве стиля
титульного листа и стиль \val{plain} для всех остальных страниц.

\subsection{Добавление листа нестандартного размера в документ}
\label{s:freesize}

\subsubsection{Общая информация}

\textbf{По состоянию на 2010-01-12 данная возможность является
экспериментальной и нуждается в тестировании и доводке. Формат команд
в будущем может быть изменен.}

Иногда бывает необходимо вставить в документ формата А4 несколько
листов формата А3 в альбомной ориентации. С помощью пакета
\pkg{eskdfreesize} можно вставить лист любого размера в любой
ориентации.

В пакете определена следующая команда:
\begin{defs}
  \cmddef{ESKDfreeSize\{\cmdarg{ширина}\}%
  \{\cmdarg{высота}\}\{\cmdarg{содержимое}\}}{}
\end{defs}

Первые~2 аргумента~--- размеры листа в единицах, понятных \LaTeX.
Третий аргумент~--- содержимое, которое необходимо разместить на
данном листе. В качестве содержимого могу выступать рисунки, таблицы,
текст и т.д.

Данная команда при вызове автоматически начинает новый лист. Если все
содержимое не поместится на один лист~--- будет начат следующий лист
такого же нестандартного размера до тех пор, пока все содержимое не
будет размещено. 

Нестандартные листы оформляются в специальном стиле \val{freesize}
(подробнее об управлении стилями см. раздел~\ref{s:pagestyles}) Данный
стиль сам подбирает размер рамки и положение элементов основной
надписи в зависимости от размеров листа.  Стиль \val{freesize} целиком
<<унаследован>> от стиля \val{formIIab} за исключением графы 32 по
причине того, что формат листа заранее не известен. По умолчанию графа
32 в данном стиле содержит только слово \cmdarg{<<Формат>>} с
последующим пробелом. Эту графу можно отключить во всем документе
(опция \opt{nocolumnxxxii}), можно удалить из стиля \val{freesize} с
помощью команды:
\begin{defs}
  \cmddef{ESKDremoveFromStyle\{\cmdarg{freesize}\}\{%
    \cmdarg{columnxxxii}\}}{,}
\end{defs}

наконец, можно вручную переопределить ее содержимое, например:
\begin{defs}
  \cmddef{renewcommand\textbackslash 
    ESKDtheColumnXXXIIfree\{\cmdarg{Формат А3}\}}{}
\end{defs}

Переопределять ее необходимо перед вызовом \cmd{ESKDfreeSize}.

Для удобства, а так же в качестве примеров в конце
\pkg{eskdfreesize.sty} были определены следующие команды:
\begin{defs}
  \cmddef{ESKDlandscapeAIII}{}
  \cmddef{ESKDportraitAIII}{}
  \cmddef{ESKDlandscapeAII}{}
  \cmddef{ESKDportraitAII}{}
  \cmddef{ESKDlandscapeAI}{}
  \cmddef{ESKDportraitAI}{}
  \cmddef{ESKDlandscapeAZERO}{}
  \cmddef{ESKDportraitAZERO}{}
  \cmddef{ESKDlandscapeBanner}{}
\end{defs}

Назначение команд понятно из их названий, все они принимают один
аргумент~--- содержимое. 
  
\subsubsection{Дополнительные возможности}  

Для удобства последующей автоматической сортировки страниц разного
формата по разным файлам (например, с помощью утилиты \pkg{pdftk})
факт вставки нестандартных листов фиксируется в \opt{jobname.log} в
следующем формате:
\small
  \begin{verbatim}
    eskdfreesize: sheet(s) with size (W x H) inserted at page(s) [RANGE]\end{verbatim}
\normalsize
\begin{ESKDexplanation} 
\item [где ] W --- ширина;
\item H --- высота;
\item RANGE --- диапазон номеров страниц в виде \mbox{[$N_1$-$N_2$]},
  либо [$N_1$], если страница только одна.
\end{ESKDexplanation}

Например, для альбомного листа А3, попавшего на 20-ю страницу,
запись будет следующей:
\small
  \begin{verbatim}
    eskdfreesize: sheet(s) with size (420mm x 297mm) inserted at page(s) [20]\end{verbatim}
\normalsize

\subsubsection{Ограничения и недостатки}

В первую очередь, ограничения по размеру.  У Adobe Acrobat ограничение
на размер листа 200х200in, у pdflatex~--- 16383x16383pt.

Следующее ограничение связано с используемым интерпретатором:
изменение размера возможно только при прямой компиляции в *.pdf c
помощью pdflatex. Не смотря на то, что latex тоже сможет обработать
такой документ, все листы, размер которых больше заданного в
\cmd{documentclass} будут обрезаны.

\subsection{Поддрежка ДСТУ (Украина)}

В некоторых случаях требования ДСТУ отличаются от ГОСТов. Для
поддержки ДСТУ существует пакет \pkg{eskddstu}, который выполняет
все необходимые настройки.

На данный момент пакет \pkg{eskddstu} выполняет только настройку
окружения \env{ESKDexplanation} в соответствии с
\mbox{ДСТУ}\ \mbox{3008-95}.

\newpage
\section{Благодарности}
\sloppy
Спасибо всем участникам почтовой рассылки \mbox{CyrTeX-ru}
\mbox{[\url{https://info.vsu.ru/Lists/CyrTeX-ru/List.html}]}, которые
помогают мне решать вопросы, связанные с разработкой \pkg{eskdx}.

\fussy

Спасибо всем, кто присылал свои предложения и сообщал об ошибках.

Отдельная благодарность Александру Роженко за его коллекцию пакетов
\pkg{ncctools}, откуда я позаимствовал некоторые принципы.

Спасибо Владимиру Пылинскому за его вклад (\pkg{eskdfreesize}, и
другие улучшения).

А также спасибо Киевскому техникуму электронных приборов, который
побудил меня на разработку коллекции \pkg{eskdx}. Если бы не курсовые
и дипломная, я бы этого не делал.


\ESKDappendix{справочное}{Расположение полей титульного листа и
  листа утверждения}
\label{a:titlelayout}
\begin{figure}[h]
\begin{center}
\includegraphics[scale=0.6]{img-title}
\end{center}
\end{figure}

\ESKDappendix{справочное}{Расположение граф основной надписи}
\label{a:columnlayout}
\begin{figure}[h]
\begin{center}
\includegraphics[scale=0.6]{img-form1}
\end{center}
\caption{Основная надпись для чертежей и схем. Форма~1}
\end{figure}

\begin{figure}[h]
\begin{center}
\includegraphics[scale=0.6]{img-form2}
\end{center}
\caption{Основная надпись для текстовых конструкторских документов
(первый или заглавный лист). Форма~2}
\end{figure}

\begin{figure}[h]
\begin{center}
\includegraphics[scale=0.6]{img-form2a}
\end{center}
\caption{Основная надпись для чертежей (схем)
и текстовых конструкторских документов (последующие листы). Форма~2а}
\end{figure}

\end{document}

